% !TEX program = xelatex
% !TEX encoding = UTF-8 Unicode
%
% 大连海事大学博士学位论文开题报告 — LaTeX 模板
% 源文件：《博士研究生开题、中期、预答辩、答辩相关表格（2022年版）》之
%        「1博士研究生学位论文开题报告.doc」
% 引擎：XeLaTeX + ctexart + biblatex (gb7714-2015)
% 使用说明见 README.md
%
% 字体规格（从原 .doc 二进制提取）：
%   封面标题   宋体加粗 一号(26pt)
%   封面信息   宋体加粗 三号(16pt)
%   节标题     黑体     小四(12pt)
%   小节标题   宋体     小四(12pt)
%   正文       宋体     五号(10.5pt)
%   提示语     楷体     五号(10.5pt)
%   评分表     宋体     五号(10.5pt)
%   每节外框   1.5px solid #BFBFBF 表格

\documentclass[UTF8, a4paper, zihao=-4, fontset=mac]{ctexart}

%==================== 宏包 ====================
\usepackage[
  top=2.5cm,
  bottom=2.5cm,
  left=3.0cm,     % 左侧留出装订余量
  right=2.5cm,
  headheight=14pt,
  headsep=0.5cm,
  footskip=1.0cm
]{geometry}

% 数学
\usepackage{amsmath}
\usepackage{amssymb}
\usepackage{mathtools}

% 表格
\usepackage{array}
\usepackage{booktabs}
\usepackage{multirow}
\usepackage{tabularx}
\usepackage{makecell}

% 图形
\usepackage{graphicx}
\usepackage{float}

% 列表
\usepackage{enumitem}

% 颜色
\usepackage{xcolor}

% 可跨页带边框框体（还原 Word 模板每节的大表格外框）
\usepackage[most]{tcolorbox}

% 页眉页脚
\usepackage{fancyhdr}

% 参考文献：GB/T 7714-2015（国标，符合学校撰写规范）
\usepackage[backend=biber, style=gb7714-2015]{biblatex}
\addbibresource{references.bib}

% 超链接（最后加载）
\usepackage[
  colorlinks=true,
  linkcolor=black,
  citecolor=blue!50!black,
  urlcolor=blue!60!black,
  bookmarksnumbered=true,
  pdfstartview=FitH,
  pdftitle={博士学位论文开题报告},
  pdfauthor={}
]{hyperref}

%==================== 排版参数 ====================
\linespread{1.2}
\setlength{\parindent}{2em}
\setlength{\parskip}{10pt plus 2pt minus 1pt}

% macOS 宋体无独立粗体字重，ctex 默认将 \songti\bfseries 映射到黑体；
% 此处覆盖为合成粗体（加粗笔画），使"加粗"表现为视觉粗体而非字体切换
\setCJKfamilyfont{zhsong}{Songti SC}[AutoFakeBold=3]

% 英文字体：Times New Roman
\setmainfont{Times New Roman}

\pagestyle{fancy}
\fancyhf{}
\fancyfoot[C]{\thepage}
\renewcommand{\headrulewidth}{0pt}

%==================== 框体与标题命令 ====================

% 每节的大框线表格（还原 Word 模板：1.5px solid #BFBFBF）
% black!25 = rgb(191,191,191) = #BFBFBF
\newtcolorbox{formbox}{%
  breakable,
  enhanced,
  colback=white,
  colframe=black!25,
  boxrule=1pt,
  arc=0pt,
  outer arc=0pt,
  left=10pt, right=10pt,
  top=10pt, bottom=10pt,
  before skip=4pt, after skip=10pt,
  fontupper=\zihao{5}\songti,
  before upper={\setlength{\parskip}{10pt plus 2pt minus 1pt}\setlength{\parindent}{2em}},
}

% 整页填写框：\clearpage 后自动撑满本页剩余高度，
% 框内留出大片空白供打印后手写（用于第 8、9、12 节意见/评议栏）
% 注意：height fill 与 breakable 互斥，此框不可跨页
\newtcolorbox{formboxfull}{%
  enhanced,
  height fill,
  valign=top,
  colback=white,
  colframe=black!25,
  boxrule=1pt,
  arc=0pt,
  outer arc=0pt,
  left=10pt, right=10pt,
  top=10pt, bottom=10pt,
  before skip=4pt, after skip=10pt,
  fontupper=\zihao{5}\songti,
  before upper={\setlength{\parskip}{10pt plus 2pt minus 1pt}\setlength{\parindent}{2em}},
}

% 节标题：黑体 小四（节号已硬编码，与学校模板一致）
\newcommand{\formheading}[1]{%
  \noindent{\heiti #1}\par\nobreak\vspace{2pt}%
}

% 小节标题：黑体 小四
\newcommand{\formsubheading}[1]{%
  \vspace{4pt}\noindent{\zihao{-4}\heiti #1}\par\nobreak\vspace{2pt}%
}

% 子小节标题（如 1.1、1.2）：宋体五号加粗
\newcommand{\formsubsubheading}[1]{%
  \vspace{4pt}\noindent{\zihao{5}\songti\bfseries #1}\par\nobreak\vspace{2pt}%
}

% 题目提示语：楷体 五号 灰色（原文括号中的填写说明）
\newcommand{\topicnote}[1]{%
  {\zihao{5}\kaishu\color{gray!50!black}#1\par\vspace{6pt}}%
}

% 空白填写区占位符（已禁用，不再产生空白；TODO 注释保留作参考）
\newcommand{\writehere}[1][6cm]{}

% 表单项标签（签名/勾选行）：黑体 五号，与正文宋体区分
\newcommand{\formlabel}[1]{{\heiti #1}}

%==================== 个人信息（提交前务必填写）====================
\newcommand{\studentName}{姓名}
\newcommand{\studentID}{学号}
\newcommand{\schoolName}{学院名称}          % TODO: 核对学院全称
\newcommand{\discipline}{学科专业}          % TODO: 核对学科名称
\newcommand{\advisorName}{指导教师}         % TODO: 导师姓名
\newcommand{\thesisTitle}{论文题目}         % TODO: 填写论文题目
\newcommand{\proposalDate}{202X 年 X 月 X 日}  % TODO: 填写开题日期
%==================================================================

\begin{document}

%==================================================
%  封面（宋体加粗，字号按原模板）
%==================================================
\begin{titlepage}
\linespread{1}\selectfont  % 封面用单倍行距，不受全局 1.2 影响
\begin{center}
\vspace*{2cm}

{\zihao{1}\songti\bfseries 大连海事大学博士研究生}\\[16pt]
{\zihao{1}\songti\bfseries 学位论文开题报告}
\end{center}

\vspace{3.5cm}

% 信息表：按原 Word 模板缩进居中，值列用 p{} 允许长题目换行
\noindent\hspace*{2.5cm}%
\renewcommand{\arraystretch}{1.8}
{\zihao{3}\songti\bfseries
\begin{tabular}{r@{：}p{8.2cm}}
  \makebox[4em][s]{姓\hfill 名} & \studentName \\
  \makebox[4em][s]{学\hfill 号} & \studentID \\
  \makebox[4em][s]{学\hfill 院} & \schoolName \\
  \makebox[4em][s]{学\hfill 科} & \discipline \\
  \makebox[4em][s]{指导教师}     & \advisorName \\
  \makebox[4em][s]{论文题目}     & \thesisTitle \\
  \makebox[4em][s]{开题时间}     & \proposalDate \\
\end{tabular}}

\vfill

\begin{center}
{\zihao{3}\songti\bfseries 大连海事大学研究生院制}\\[6pt]
{\zihao{-4}\songti\bfseries 2021年4月}

\vspace*{2cm}
\end{center}
\end{titlepage}

%==================================================
%  开题报告说明（不在框线内）
%==================================================
\begin{center}
{\zihao{3}\heiti 博士研究生学位论文开题报告说明}
\end{center}

\vspace{6pt}

{\fangsong
博士学位论文开题报告是做好学位论文的基础，为了完善博士学位过程管理体系，提高博士学位质量，原则上应于第3学期结束前完成学位论文的开题工作。

一、开题报告应在导师（组）指导下，由研究生本人完成。开题报告中文献综述应分类清楚，归纳合理，科学技术问题凝练准确，逻辑清晰，语句通顺，图表规范。字数不少于10000字（其中文献综述6000字左右）；近五年的参考文献不少于总数的1/3；参考文献的引用格式需符合《大连海事大学博士、硕士研究生学位论文撰写规范》的相关规定。严禁伪造和抄袭开题报告。

二、博士生在查阅大量国内外文献资料的基础上，填写完成《博士研究生学位论文开题报告》，经导师审核同意后，应在本学科或相关学科范围内举行公开的学位论文开题报告会。

三、开题报告会由博士生招生学科组织具体实施。开题报告考核小组负责填写结论性的审查意见，并将结果和相关材料留学院备案。开题报告考核小组由3～5名具有高级技术职称的相关学科专家组成，设组长和秘书各1名，组长必须由博士生导师担任，博士生导师人数不少于2/3，鼓励有条件的学科聘请1名外单位博士生导师作为成员。申请人的导师不能作为考核小组成员。同时，可邀请本专业领域的教师和学生参加，听取多方面意见。

四、《博士研究生学位论文开题报告》必须采用A4纸双面打印，左侧装订成册，各栏空格不够时，请自行加页。本表可在研究生院网站下载。
}

\clearpage

%==================================================
%  1、学位论文选题背景、研究目的和意义
%==================================================
\formheading{1、学位论文选题背景、研究目的和意义}
\begin{formbox}

\formsubheading{（1）选题的科学依据，提出将要研究的科学或技术问题}

% TODO: 阐述选题背景，建议按以下结构展开：
%   1.1 宏观背景（行业/领域面临的形势与挑战）
%   1.2 具体问题（聚焦到本研究关注的核心矛盾）
%   1.3 现实需求（为什么需要研究这个问题）
% 可使用 \formsubsubheading{1.1 标题} 创建子标题

\writehere[6cm]

\formsubheading{（2）研究目的，选题的理论意义或实际应用价值}

% TODO: 阐明研究目的、理论意义和实际应用价值

\writehere[6cm]
\end{formbox}

\clearpage

%==================================================
%  2、文献综述（不少于 6000 字）
%==================================================
\formheading{2、文献综述（不少于6000字）}
\begin{formbox}

\topicnote{围绕研究方向，通过有效的文献检索和大量文献阅读，运用科学分类和归纳方法，总结分析国内外的研究现状与发展趋势，凝练存在的科学或技术问题，列出主要参考文献（近五年的参考文献不少于总数的1/3）。}

% TODO: 按研究方向分 2~4 个主题展开文献综述，每个主题用 \formsubsubheading{} 标题
% 建议：
%   2.1 核心概念与理论基础
%   2.2 主要研究方法与进展
%   2.3 研究现状评述与不足（引出本研究的切入点）
% 正文中用 \cite{citekey} 或 \parencite{citekey} 引用参考文献

\formsubsubheading{2.1 主题一}

% TODO: 文献综述内容

\formsubsubheading{2.2 主题二}

% TODO: 文献综述内容

\formsubsubheading{2.3 研究现状评述}

% TODO: 总结现有研究的不足，引出本研究的科学/技术问题

\vspace{1em}
\noindent{\zihao{-4}\heiti 参考文献}
\vspace{4pt}

{\zihao{-5}\printbibliography[heading=none]}

\end{formbox}

\clearpage

%==================================================
%  3、研究内容
%==================================================
\formheading{3、研究内容}
\begin{formbox}

\formsubheading{（1）主要研究内容}

% TODO: 列出论文的主要研究内容（按篇章/模块展开）
% 建议按逻辑递进列出 3~5 个研究部分，每个部分简要说明研究内容

\writehere[7cm]

\formsubheading{（2）拟解决的关键问题}

% TODO: 凝练 2~4 个关键科学/技术问题

\writehere[5cm]
\end{formbox}

\clearpage

%==================================================
%  4、主要创新点
%==================================================
\formheading{4、主要创新点}
\begin{formbox}

% TODO: 列出 2~4 个主要创新点，每个创新点建议包含：
%   - 现有研究的不足
%   - 本研究的创新之处
%   - 创新的理论/实际意义

% 示例格式：
% （1）创新点标题。具体阐述……
% （2）创新点标题。具体阐述……

\end{formbox}

\clearpage

%==================================================
%  5、研究方案
%==================================================
\formheading{5、研究方案}
\begin{formbox}

\formsubheading{（1）拟采取的研究方法}

% TODO: 阐述拟采用的研究方法与技术手段

\formsubheading{（2）技术路线}

% TODO: 描述技术路线，可配合技术路线图
% 插图示例：
% \begin{figure}[H]
%   \centering
%   \includegraphics[width=0.85\textwidth]{figures/tech-roadmap}
%   \caption{研究技术路线图}
%   \label{fig:tech-roadmap}
% \end{figure}

\formsubheading{（3）可行性分析}

% TODO: 从理论、技术、数据、研究基础等方面分析可行性

\formsubheading{（4）预设研究中可能遇到的难点，并提出解决的方法和措施}

% TODO: 列出潜在难点及对应的解决方案
% 建议格式：
%   {\bfseries 1）难点名称}\par
%   {\bfseries 难点：}……\par
%   {\bfseries 解决方法：}……\par

\writehere[3cm]
\end{formbox}

\clearpage

%==================================================
%  6、研究基础
%==================================================
\formheading{6、研究基础}
\begin{formbox}

\formsubheading{（1）工作条件评估}

\topicnote{（科学评估工作条件，包括自己的专业知识积累、现有的科研和实验条件，估算论文工作量，协助导师具体指导的人员情况（导师组）等）}

% TODO: 从以下方面阐述工作条件：
%   1) 专业知识积累
%   2) 科研与实验条件
%   3) 论文工作量估算
%   4) 导师组及协助指导人员情况

\formsubheading{（2）经费预算}
\topicnote{（包含经费来源、开支预算）}

% TODO: 填写经费来源说明
\formsubsubheading{1）经费来源}

% TODO: 说明经费来源

\formsubsubheading{2）开支预算}

% 开支预算表格模板（按需修改项目和金额）
\begin{center}
\renewcommand{\arraystretch}{1.5}
\setlength{\tabcolsep}{6pt}
\begin{tabularx}{\textwidth}{|>{\raggedright\arraybackslash}p{4cm}|>{\raggedright\arraybackslash}X|>{\centering\arraybackslash}p{2.2cm}|}
\hline
\textbf{开支项目} & \textbf{主要用途} & \textbf{金额（元）} \\
\hline
调研差旅费 & & \\
\hline
数据收集与处理费 & & \\
\hline
 & & \\
\hline
\textbf{总计} & & \textbf{} \\
\hline
\end{tabularx}
\end{center}

\writehere[5cm]
\end{formbox}

\clearpage

%==================================================
%  7、学位论文工作计划（本节本身就是表格，不需要 formbox 外框）
%==================================================
\formheading{7、学位论文工作计划}

\vspace{0.3cm}

\renewcommand{\arraystretch}{1.4}
\setlength{\tabcolsep}{4pt}
\noindent
\begin{tabularx}{\textwidth}{|>{\centering\arraybackslash}p{3cm}|>{\centering\arraybackslash}X|>{\centering\arraybackslash}X|>{\centering\arraybackslash}X|}
\hline
\textbf{起讫日期} & \textbf{主要研究内容} & \textbf{预期目标} & \textbf{预期成果及形式} \\
\hline
 & & & \\
\hline
 & & & \\
\hline
 & & & \\
\hline
 & & & \\
\hline
 & & & \\
\hline
\end{tabularx}

\vspace{0.3cm}

{\zihao{5}\color{gray!50!black}注：表格行数不够时可自行添加。}

%==================================================
%  8、导师的意见（单独一页）
%==================================================
\clearpage
\formheading{8、导师的意见}
\begin{formboxfull}

\topicnote{（导师须明确该题目是否属于本学科研究领域范畴，如属于交叉学科，则须明确交叉领域与本学科研究领域的具体关系；需要对选题依据、研究思路、研究内容和技术路线的科学性、可行性、先进性及创新性进行评价）}

\vfill
\vspace{10cm}

\noindent\formlabel{导师是否能够提供学位论文工作所需的实验条件和研究经费保障：}\underline{\hspace{6em}}


\noindent\formlabel{导师是否同意该博士生举行开题报告会：}\underline{\hspace{10em}}


\noindent\formlabel{导师签字：}\underline

\vspace{1.2cm}

{\hspace{5em}}\hfill \proposalDate
\end{formboxfull}

\clearpage
%==================================================
%  9、点长意见（如学院设立点长）
%==================================================
\formheading{9、博士学位授权点点长对开题报告的意见}
\begin{formboxfull}

\topicnote{（如学院设立点长，请签署此项。重点指出文献综述是否进行了充分的凝练和概括，能否反映相关国内外的研究现状和概况，论文选择研究的科学或者技术问题是否清晰明确、是否具有学术价值或应用潜力，研究内容是否属于本学科研究领域或范畴，需要对选题依据、研究思路、研究内容和技术路线的科学性、可行性、先进性及创新性进行评价）}

\vfill

\vspace{16cm}

\noindent\formlabel{点长是否同意该博士生举行开题报告会：}\underline{\hspace{10em}}
\vspace{1.2cm}

\noindent\formlabel{点长签字：}\underline{\hspace{5em}}\hfill \proposalDate
\end{formboxfull}

\clearpage
%==================================================
%  10、评审小组组成
%==================================================
\formheading{10、评审小组组成}

\topicnote{（评审小组成员组成由学科带头人审核）}

\vspace{0.3cm}

{\zihao{5}
\renewcommand{\arraystretch}{1.8}
\setlength{\tabcolsep}{3pt}
\noindent
\begin{tabular}{|>{\centering\arraybackslash}p{1.6cm}|>{\centering\arraybackslash}p{2.3cm}|>{\centering\arraybackslash}p{2cm}|>{\centering\arraybackslash}p{2.4cm}|>{\centering\arraybackslash}p{3.4cm}|>{\centering\arraybackslash}p{2cm}|}
\hline
\textbf{组成} & \textbf{姓名} & \textbf{职称} & \textbf{学科\newline（专业）} & \textbf{工作单位\newline（本校专家不填）} & \textbf{签字} \\
\hline
组长 & & & & & \\
\hline
\multirow{4}{*}{成员} & & & & & \\
\cline{2-6}
 & & & & & \\
\cline{2-6}
 & & & & & \\
\cline{2-6}
 & & & & & \\
\hline
\end{tabular}}

\vspace{0.8cm}

\hfill\formlabel{审核人签字：}\underline{\hspace{5em}}

%==================================================
%  11、开题报告会记录
%==================================================
\formheading{11、开题报告会记录}
\begin{formbox}

\topicnote{（记录评审专家的质疑问题与研究生的回答要点，以及专家对选题的具体修改意见）}

\writehere[100cm]

\vspace{10cm}

\noindent\formlabel{记录人签字：}\underline{\hspace{5em}} \hfill \proposalDate
\end{formbox}

%==================================================
%  12、评审专家小组评议（单独一页）
%==================================================
\clearpage
\formheading{12、评审专家小组对开题报告的评议}
\begin{formboxfull}

\formsubheading{（1）对选题依据、研究思路、研究内容和技术路线的科学性、可行性、先进性及创新性的评价}

\vfill

\vspace{4cm}

\formsubheading{（2）存在的主要问题和改进措施}

\vfill

\vspace{4cm}

\formsubheading{（3）评定结果}

\vspace{0.2cm}

\hspace{2em}$\square$~优秀（90$\sim$100分）\qquad $\square$~良好（80$\sim$89分）

\vspace{0.2cm}

\hspace{2em}$\square$~合格（60$\sim$79分）\qquad $\square$~不合格（60分以下）

\vspace{0.4cm}

\formsubheading{（4）评定结论}

\vspace{0.2cm}

\hspace{2em}$\square$~通过 \qquad $\square$~不通过，需要重新开题

\vspace{0.8cm}

\noindent 评审专家小组

\vspace{0.3cm}

\noindent\formlabel{组长签字：}\underline{\hspace{5em}}

\vspace{0.3cm}

\noindent\formlabel{成员签字：}\underline{\hspace{12em}}

\vspace{0.3cm}

\noindent\formlabel{秘书签字：}\underline{\hspace{5em}} \hfill \proposalDate
\end{formboxfull}

%==================================================
%  附录：评分表（五号字，独立表格，不在 formbox 内）
%==================================================
\clearpage
\begin{center}
{\zihao{3}\fangsong 附：评审小组专家对博士研究生学位论文开题报告评分表}
\end{center}

\vspace{0.4cm}

{\zihao{5}\fangsong
\renewcommand{\arraystretch}{1.2}
\setlength{\tabcolsep}{4pt}
\noindent
\begin{tabularx}{\textwidth}{|p{2.4cm}|c|X|c|}
\hline
\textbf{评审项目} & \textbf{权重} & \textbf{评分标准} & \makecell{\textbf{得分}\\\textbf{(百分制)}} \\
\hline
\makecell[l]{一、选题依据\\(A)} & 20\% &
\textbf{80$\sim$100分：}选题有很强的理论意义、实用价值，并预期将获得重大的经济效益和社会效益。
\textbf{60$\sim$80分：}选题有较强的理论意义、实用价值，并预期获得较大的经济效益和社会效益。
\textbf{60分以下：}选题缺乏理论意义和实用价值。 & \\
\hline
\makecell[l]{二、文献综述\\(B)} & 10\% &
\textbf{80$\sim$100分：}查阅大量文献，报告内容能全面阐述该研究方向的现状和发展动态。
\textbf{60$\sim$80分：}查阅一定量文献，报告内容基本跟踪该研究方向的现状和发展动态。
\textbf{60分以下：}查阅少量文献，综述一般，未达到上述标准。 & \\
\hline
\makecell[l]{三、研究内容\\和技术路线\\(C)} & 20\% &
\textbf{90$\sim$100分：}能够提出重要的科学问题，研究思路清晰，研究内容充实先进，技术路线科学、合理、有可行性。
\textbf{80$\sim$90分：}能够提出比较重要的科学问题，研究思路较清晰，研究内容较充实，技术路线合理、有一定可行性。
\textbf{60$\sim$80分：}能够提出科学问题，研究思路一般，研究内容尚可，技术路线有一定可行性，但有不确定可能。
\textbf{60分以下：}未提出明确的科学问题，研究思路不清晰，研究内容不充分，技术路线缺乏可行性。 & \\
\hline
\makecell[l]{四、创新性\\(D)} & 30\% &
\textbf{80$\sim$100分：}研究课题属本学科发展方向并居前沿位置，具有很强的创新性。
\textbf{60$\sim$80分：}研究课题属本学科的发展方向，有自己独特的思考、具有一定的创新性。
\textbf{60分以下：}研究课题归属不清，主要创新点不明显。 & \\
\hline
\makecell[l]{五、文字表达\\(E)} & 10\% &
\textbf{80$\sim$100分：}条理清晰，分析严谨，文笔流畅。
\textbf{60$\sim$80分：}条理较好，层次分明，文笔较流畅。
\textbf{60分以下：}写作能力较差。 & \\
\hline
\makecell[l]{六、口头报告\\(F)} & 10\% &
\textbf{80$\sim$100分：}基本概念明确，思路清晰、逻辑性强、表达清楚。
\textbf{60$\sim$80分：}基本概念明确，思路基本清晰、层次较分明、表达较清楚。
\textbf{60分以下：}基本概念和思路不清晰、逻辑混乱、表达不清楚。 & \\
\hline
\textbf{总分} & & \multicolumn{2}{l|}{总分 $= 0.2A + 0.1B + 0.2C + 0.3D + 0.1E + 0.1F$} \\
\hline
\end{tabularx}

\vspace{0.5cm}

\textbf{注：}此表可复印，每位评审专家在六项指标每一栏的最后一列内打分（百分制），所有评审表均装订入开题报告。
}

\end{document}
